\documentclass[main.tex]{subfiles}
\begin{document}

\section{Capital allocation across strategies}\label{sec:alloc}

\subsection{Motivation and setup}
Most strategies in this study commit only part of their capital. The option sleeves cap risk at 0.5--1\% of equity per position, and the event sleeves hold positions on a fraction of days. A combined book therefore needs an explicit rule for capital that is not deployed. We park it in the Secured Overnight Financing Rate (SOFR). We then ask whether a portfolio of sleeves, with weights fixed from 2022--2023 data alone, would have added return over SOFR in 2024.

Let $x_{i,t}$ be the daily excess return of sleeve $i$ over cash. These are the series the C++ engines write, one per pre-registered primary; the 13 sleeves are listed in Table~\ref{tab:sleeves}. Let $c_t$ be the SOFR accrual from session $t-1$ to $t$, using the fixing of session $t-1$ and an actual/360 day count. The New York Fed publishes the series; we downloaded it from its public API (770 fixings, 2021-12-01 to 2024-12-31). The portfolio total return is
\begin{equation}
  r^{p}_{t} \;=\; c_t \;+\; \sum_{i=1}^{13} w_i\, x_{i,t}, \qquad w_i \ge 0,\quad \sum_i w_i \le 1 .
  \label{eq:port}
\end{equation}
A weight $w_i$ is the share of capital run through sleeve $i$. Inside a sleeve, idle cash already earns the cash rate, so \eqref{eq:port} credits SOFR on all undeployed capital. The sleeves' own excess returns are measured against the 1-month T-bill. We treat the small T-bill/SOFR basis as negligible, which is an approximation.

\subsection{Allocation rules}
All rules are long-only. Covariances use the Ledoit--Wolf shrinkage estimator toward a scaled identity \citep{ledoit2004}.
\begin{enumerate}[leftmargin=*]
  \item \textbf{1/N}: equal capital to every sleeve. This is the benchmark that estimated rules often fail to beat out of sample \citep{demiguel2009}.
  \item \textbf{Inverse volatility}: $w_i \propto 1/\hat\sigma_i$, normalized to sum to one.
  \item \textbf{Minimum variance}: $\min_w w^\top \hat\Sigma w$ subject to $\sum_i w_i = 1$.
  \item \textbf{Markowitz (Bayes--Stein, 25\% cap)}: maximize $w^\top \tilde\mu / \sqrt{w^\top\hat\Sigma w}$ subject to $0 \le w_i \le 0.25$ and $\sum_i w_i = 1$ (the Sharpe ratio is scale-free, so the budget is fixed inside the optimizer) \citep{markowitz1952}. Here $\tilde\mu$ shrinks sample means toward the minimum-variance portfolio mean with Jorion's intensity \citep{jorion1986}.
  \item \textbf{Development-screened 1/N}: equal weight on sleeves whose estimation-window Newey--West $t$ exceeds 1.
\end{enumerate}

\textbf{Schedules.} Weights are estimated on 2022-01-03 to 2023-12-29 and applied two ways:
\begin{itemize}[leftmargin=*]
  \item \textbf{Static}: estimated once and held through 2024.
  \item \textbf{Walk-forward}: re-estimated at each 2024 month start on an expanding window.
\end{itemize}

\textbf{Volatility target.} We also scale the static 1/N and Markowitz books to a 5\% in-sample volatility, with leverage capped at 3, following the logic of volatility targeting \citep{moreira2017}. Leverage multiplies excess returns, and financing is implicitly at SOFR.

\subsection{Results}
Table~\ref{tab:alloc} reports each rule. In 2024, holding SOFR alone returned 5.40\% (mean annualized accrual 5.26\%).

\begin{table}[ht]\centering\small
\caption{Allocation results. Excess and Sharpe are relative to SOFR. Weights never use 2024 data, except that the walk-forward schedule uses 2024 data up to each rebalance.}\label{tab:alloc}
\begin{tabular}{llrrrrr}
\toprule
Rule & Schedule & Total & Excess & Vol & Sharpe & Max DD \\
\midrule
\multicolumn{7}{l}{\emph{2022--2023, in-sample}} \\
1/N & static & 3.93\% & 0.50\% & 2.03\% & 0.25 & $-1.34\%$ \\
Inverse vol & static & 3.51\% & 0.08\% & 0.76\% & 0.11 & $-0.42\%$ \\
Minimum variance & static & 3.33\% & $-0.10\%$ & 0.42\% & $-0.23$ & $-0.30\%$ \\
Markowitz (B--S, cap 25\%) & static & 6.22\% & 2.67\% & 1.65\% & 1.62 & $-0.86\%$ \\
Dev-screened 1/N & static & 7.42\% & 3.84\% & 3.10\% & 1.24 & $-1.93\%$ \\
\midrule
\multicolumn{7}{l}{\emph{2024, out-of-sample}} \\
1/N & static & 9.91\% & 4.21\% & 2.10\% & 2.00 & $-0.89\%$ \\
Inverse vol & static & 5.91\% & 0.49\% & 1.01\% & 0.48 & $-0.45\%$ \\
Inverse vol & walk-forward & 5.99\% & 0.57\% & 1.00\% & 0.56 & $-0.43\%$ \\
Minimum variance & static & 5.16\% & $-0.23\%$ & 0.55\% & $-0.41$ & $-0.23\%$ \\
Minimum variance & walk-forward & 5.12\% & $-0.27\%$ & 0.51\% & $-0.53$ & $-0.22\%$ \\
Markowitz (B--S, cap 25\%) & static & 9.31\% & 3.67\% & 1.94\% & 1.89 & $-0.74\%$ \\
Markowitz (B--S, cap 25\%) & walk-forward & 10.43\% & 4.69\% & 2.27\% & 2.07 & $-0.89\%$ \\
Dev-screened 1/N & static & 10.05\% & 4.35\% & 2.67\% & 1.63 & $-1.13\%$ \\
Dev-screened 1/N & walk-forward & 14.01\% & 7.91\% & 3.54\% & 2.23 & $-1.13\%$ \\
1/N, 5\% vol target ($\times 2.46$) & static & 16.75\% & 10.37\% & 5.18\% & 2.00 & $-2.29\%$ \\
Markowitz, 5\% vol target ($\times 3.00$) & static & 17.45\% & 11.00\% & 5.82\% & 1.89 & $-2.61\%$ \\
\bottomrule
\end{tabular}
\end{table}


\Fig{alloc}{fig18_alloc_weights.png}{Static weights estimated on 2022--2023 for each allocation rule. The dashed line marks the 25\% Markowitz cap.}{fig:alloc_w}
\Fig{alloc}{fig19_alloc_corr.png}{Correlations of daily sleeve returns, 2022--2023.}{fig:alloc_corr}
\Fig{alloc}{fig20_alloc_2024.png}{Cumulative 2024 total return of selected allocation rules, including SOFR on idle capital.}{fig:alloc_2024}

Figures~\ref{fig:alloc_w}--\ref{fig:alloc_2024} show the weights, the sleeve correlations (median off-diagonal $-0.01$, largest absolute value 0.25) and the 2024 paths.

\textbf{Reading the table.}
\begin{itemize}[leftmargin=*]
  \item \textbf{Minimum variance.} It loads on the sleeves with the lowest volatility: the small-sized overnight and pre-earnings straddle books (H2, H5), which carry 0.5--0.6\% annual volatility at their risk caps (2022--2023). It returns roughly SOFR. Low volatility here reflects small position sizes, not diversification.
  \item \textbf{Markowitz with shrunk means.} It avoids that trap. It puts its weight on the sleeves with positive estimation-window means. H6 and S1 reach the 25\% cap; S2, S3 and S9 receive 13--15\%; H5, S8 and S10 receive at most 4\%. It earns 3.7--4.7\% over SOFR in 2024.
  \item \textbf{1/N.} It earns 4.2\% over SOFR. That return comes largely from sleeves that were weak in 2022--2023 and strong in 2024 (H3 overnight equity, S8): equal weighting kept them in the book.
  \item \textbf{Uncertainty.} No 2024 Sharpe ratio in Table~\ref{tab:alloc} rests on more than 252 daily observations. The standard error of a one-year Sharpe estimate is about 1, so the differences between the 1/N, Markowitz and screened rules are not statistically distinguishable.
\end{itemize}

\subsection{Why this 2024 test is weaker than it looks}
Three caveats limit what Table~\ref{tab:alloc} can support.
\begin{enumerate}[leftmargin=*]
  \item \textbf{Known sleeve outcomes.} No weight uses 2024 returns, but the sleeves' 2024 outcomes were already known when this allocation study was designed: the wave-2 holdout and the wave-3 test both cover 2024. The set of sleeves is ``all primaries'', which limits but does not remove that knowledge. The option-flow sleeve uses the 1-day hold (the wave-2 H4 specification with realistic fees), not the 5-day hold designated in the wave-3 pre-registration; the estimation window would have chosen it as well (Sharpe 0.95 against $-0.46$).
  \item \textbf{Fees and contract sizes.} Wave-2 sleeves (H1--H6) carry the wave-2 fee assumptions, which understate single-name and ETF option fees by about \$0.43 per contract. Scaling a sleeve by $w_i$ also ignores integer contract sizes. At $w_i = 0.077$ on a \$1M account, several option sleeves would trade fractional contracts.
  \item \textbf{Leverage.} The volatility-targeted books use 2.46--3 times leverage on sleeves whose risk is dominated by rare tail days (H1, S5). Their 2024 Sharpe ratios are the unlevered ratios. Their drawdowns would scale with any tail event that 2024 did not contain.
\end{enumerate}

\subsection{Allocation proposal}
These results suggest a combined book can add return over SOFR with low volatility, but they do not establish it. We propose a staged allocation in which every stage has a pre-registered gate.

\begin{enumerate}[leftmargin=*]
  \item \textbf{Default: SOFR.} All capital not assigned to a sleeve earns SOFR, through overnight reverse repo or a government money-market fund. This is the benchmark every stage must beat after costs.
  \item \textbf{Stage 1, paper portfolio in 2025.}
    \begin{itemize}[leftmargin=*]
      \item Freeze the Bayes--Stein Markowitz rule (25\% cap, monthly walk-forward on an expanding window) and the 1/N rule as a control.
      \item Universe: the 13 sleeves, with wave-2 sleeves re-costed at the wave-3 fee schedule.
      \item Run both on 2025 data without changes.
      \item Pass gate: positive excess over SOFR, a Sharpe above that of the 1/N control, and a drawdown below 2\% unlevered.
    \end{itemize}
  \item \textbf{Stage 2, small live risk budget.} Only if Stage 1 passes, allocate to the two sleeves with replicated evidence (S3 option flow and S2 announcement drift) at a combined 5\% annualized volatility budget, with the remainder in SOFR. Each sleeve has a kill switch: stop if its live rolling 63-session Sharpe is below $-1$, or if its drawdown exceeds twice its 2022--2024 maximum.
  \item \textbf{Excluded sleeves.} The short-volatility sleeves (H1, S5) and the straddle sleeves (H2, H5, S6) have negative estimation-window means and receive zero Markowitz weight. They stay out unless a new pre-registered test supports them.
\end{enumerate}

This proposal describes a research protocol for the strategies in this study. It is not investment advice. The allocation code is \texttt{alloc/allocate.py}; weights, sleeve statistics and daily portfolio returns are in \texttt{alloc/}.

\ifSubfilesClassLoaded{\section*{References}% Bibliography.
\bigskip\noindent\emph{The reference details below should be checked against the published versions before citing.}\par\medskip
\begin{thebibliography}{99}
\setlength{\itemsep}{2pt}
\small
\bibitem[Bailey and L{\'o}pez de Prado(2014)]{bailey2014} Bailey, D. H., \& L{\'o}pez de Prado, M. (2014). The deflated Sharpe ratio. \emph{Journal of Portfolio Management}, 40(5).
\bibitem[Barber et~al.(2013)]{barber2013} Barber, B. M., De George, E. T., Lehavy, R., \& Trueman, B. (2013). The earnings announcement premium around the globe. \emph{Journal of Financial Economics}, 108(1).
\bibitem[Brandt et~al.(2008)]{brandt2008} Brandt, M. W., Kishore, R., Santa-Clara, P., \& Venkatachalam, M. (2008). Earnings announcements are full of surprises. Working paper.
\bibitem[Chan et~al.(1996)]{chan1996} Chan, L. K. C., Jegadeesh, N., \& Lakonishok, J. (1996). Momentum strategies. \emph{Journal of Finance}, 51(5).
\bibitem[Cremers and Weinbaum(2010)]{cremers2010} Cremers, M., \& Weinbaum, D. (2010). Deviations from put-call parity and stock return predictability. \emph{Journal of Financial and Quantitative Analysis}, 45(2).
\bibitem[DeMiguel et~al.(2009)]{demiguel2009} DeMiguel, V., Garlappi, L., \& Uppal, R. (2009). Optimal versus naive diversification: How inefficient is the 1/N portfolio strategy? \emph{Review of Financial Studies}, 22(5).
\bibitem[Easley et~al.(1998)]{easley1998} Easley, D., O'Hara, M., \& Srinivas, P. S. (1998). Option volume and stock prices: Evidence on where informed traders trade. \emph{Journal of Finance}, 53(2).
\bibitem[Fama and MacBeth(1973)]{famamacbeth1973} Fama, E. F., \& MacBeth, J. D. (1973). Risk, return, and equilibrium: Empirical tests. \emph{Journal of Political Economy}, 81(3).
\bibitem[Frazzini and Lamont(2007)]{frazzini2007} Frazzini, A., \& Lamont, O. A. (2007). The earnings announcement premium and trading volume. NBER Working Paper 13090.
\bibitem[Johnson and So(2012)]{johnsonso2012} Johnson, T. L., \& So, E. C. (2012). The option to stock volume ratio and future returns. \emph{Journal of Financial Economics}, 106(2).
\bibitem[Jorion(1986)]{jorion1986} Jorion, P. (1986). Bayes-Stein estimation for portfolio analysis. \emph{Journal of Financial and Quantitative Analysis}, 21(3).
\bibitem[Ledoit and Wolf(2004)]{ledoit2004} Ledoit, O., \& Wolf, M. (2004). A well-conditioned estimator for large-dimensional covariance matrices. \emph{Journal of Multivariate Analysis}, 88(2).
\bibitem[Markowitz(1952)]{markowitz1952} Markowitz, H. (1952). Portfolio selection. \emph{Journal of Finance}, 7(1).
\bibitem[Newey and West(1987)]{neweywest1987} Newey, W. K., \& West, K. D. (1987). A simple, positive semi-definite, heteroskedasticity and autocorrelation consistent covariance matrix. \emph{Econometrica}, 55(3).
\bibitem[Pan and Poteshman(2006)]{pan2006} Pan, J., \& Poteshman, A. M. (2006). The information in option volume for future stock prices. \emph{Review of Financial Studies}, 19(3).
\bibitem[Xing et~al.(2010)]{xing2010} Xing, Y., Zhang, X., \& Zhao, R. (2010). What does the individual option volatility smirk tell us about future equity returns? \emph{Journal of Financial and Quantitative Analysis}, 45(3).
\end{thebibliography}
}{}
\end{document}
